\section{Explicit Numerical Bounds and Physical Predictions}
\label{sec:numerical-bounds}
%=============================================================================

\begin{theorem}[Mass Gap Estimates (Framework Prediction)]
For the physical string tension $\sqrt{\sigma_{\text{phys}}} \approx 440$ MeV, effective string theory predicts:

\textbf{$SU(2)$}:
\[
\Delta_{SU(2)} \geq \sqrt{\frac{2\pi}{3}} \cdot \sqrt{\sigma} \approx 1.45 \cdot 440\text{ MeV} \approx 640\text{ MeV}
\]

\textbf{$SU(3)$}:
\[
\Delta_{SU(3)} \geq \sqrt{\frac{2\pi}{3}} \cdot \frac{4}{3} \cdot \sqrt{\sigma} \approx 1.67 \cdot 440\text{ MeV} \approx 735\text{ MeV}
\]

The rigorous lower bound is $\Delta \geq \frac{2}{N}\sqrt{\sigma}$. The numerical values above should be read as \\emph{order-of-magnitude framework estimates} built on effective-string heuristics; they are not a precision prediction for the lightest glueball mass. In particular, the $SU(3)$ estimate (about $0.85$ GeV for $\sqrt{\sigma}\approx 440$ MeV) is below the commonly quoted lattice value $m_{0^{++}} \approx 1.5$--$1.7$ GeV, so it should be interpreted as a conservative lower-scale estimate rather than a direct match.
\end{theorem}

\begin{theorem}[Complete Resolution Table]
~

\begin{center}
\begin{tabular}{|l|l|l|}
\hline
\textbf{Component} & \textbf{Main Proof} & \textbf{Alternative Approach} \\
\hline
1. No phase transition & Bessel-Nevanlinna (Thm~\ref{thm:bessel-su2},\ref{thm:bessel-su3}) & Framework 4 \\
2. String tension $\sigma > 0$ & GKS/Characters (Thm~\ref{thm:sigma-positive}) & Framework 1 (Vortex) \\
3. Mass gap $\Delta > 0$ & Giles-Teper (Thm~\ref{thm:giles-teper}) & Framework 2 (Spectral) \\
4. Explicit bounds & Character expansion & Framework 3 (Tropical) \\
5. Continuum limit & OS reconstruction & Uniform bounds \\
\hline
\end{tabular}
\end{center}
\end{theorem}

%=============================================================================





